\section{Limitations and Future Work}
\label{sec:rs-limitations}

Any protocol that enforces accountability structurally — rather than by policy or convention — accepts concrete costs in exchange for that guarantee. The Recursive Spawning Protocol pays three distinct costs, each defining a direct line of future work. These costs are not implementation artifacts; they are the unavoidable trade-offs of the architectural choice to make parent pointers immutable and spawn chains tamper-evident.

% ----------------------------------------------------------------------------%
\subsection{Performance Overhead}
\label{sec:rs-perf-overhead}

The three-layer validation sequence at every spawn event introduces latency that conventional forking mechanisms do not incur. A spawn under RSP must sequentially pass Purpose Layer authorization, Bounds Engine depth evaluation, and Fact Layer warrant generation before a child activates. In the AgentOS runtime, the full sequence adds \numrange{15}{40} milliseconds per spawn under nominal load; cryptographic signing and Merkle-tree hash-chaining in the Fact Layer account for approximately 60\% of this cost. A single \texttt{fork()} in a conventional runtime completes in microseconds.

This overhead is not an implementation defect — it is the structural price of the immutable parent pointer and the tamper-evident forensic ledger. High-frequency spawn patterns — per-document children in batch workloads, per-step children in generative agent loops — accumulate this latency and degrade throughput relative to naive forking. The prescribed mitigation is batched spawning: provision fewer children, each handling multiple work units, trading decomposition granularity for lower cumulative overhead. This is contemplated by the spawn necessity threshold in the Bounds Engine.

The forensic ledger's append-only Merkle-tree structure also introduces storage and I/O overhead that compounds in long-running deployments. Each spawn event requires a signed write to the primary store and at minimum one replicated write for tamper-evidence. Write-ahead log buffering and periodic checkpoint compaction are the current mitigations; ledger growth management in sustained high-throughput deployments remains an open engineering challenge.

% ----------------------------------------------------------------------------%
\subsection{Scalability Constraints}
\label{sec:rs-scalability}

RSP's static depth bound creates a structural ceiling that conflicts with workloads requiring genuine deep decomposition. The depth bound $D_{\max}$ is fixed at the root node and applies uniformly across the chain; it cannot adapt to subtree workload characteristics without a Purpose Layer directive. A chain engaged in a task requiring authentic seven-level parallelism cannot proceed beyond $D_{\max}$ without human intervention, even when every spawn is fully authorized. Human intervention introduces latency unacceptable in real-time workloads.

The inability to dynamically adjust depth is a deliberate design property: relaxing it at runtime by any machine actor permits unbounded chain growth and collapses the termination guarantee. Future work should investigate \emph{hierarchical depth budgeting}, in which the Purpose Layer pre-authorizes a depth budget for a subtree at provisioning time, enabling bounded adaptive growth without per-spawn human involvement. The budget ceiling remains non-negotiable; a machine-extendable budget is equivalent to removing the bound.

The single-ledger architecture also constrains write throughput at scale. The forensic ledger is the authoritative record for all concurrent chains; at approximately \numrange{e4}{e5} concurrent agents under sustained spawn activity, ledger sharding by session namespace — with cross-shard audit via a sparse Merkle tree — becomes necessary. The hash-chaining property is locality-agnostic and compatible with this approach.

A third constraint is the human anchor as single escalation terminus. Simultaneous escalations from multiple chains converge on the same human principal, who becomes a sequential bottleneck. A Purpose Layer proxy issuing \texttt{ACK} directives for pre-authorized escalation categories would reduce per-escalation involvement while preserving human accountability for \texttt{RESOLVE} and \texttt{REJECT}. The proxy must be strictly limited to pre-authorized categories.

% ----------------------------------------------------------------------------%
\subsection{Compromised Purpose Layer}
\label{sec:rs-purpose-compromise}

RSP's guarantees are contingent on the human anchor's continued availability, competence, and good faith. If unavailable, the chain stalls at $D_{\text{ESCALATE}}$. If the human issues a directive without adequate diagnostic context — cognitive load, miscommunication, adversarial manipulation — the protocol has no mechanism to contest or qualify it. If the human anchor is compromised through credential theft or coercion, the entire chain above that anchor is subverted. None of RSP's correctness properties hold against a compromised Purpose Layer, because the Purpose Layer is the architectural root of all three.

This is the irreducible trade-off of any human-in-the-loop accountability architecture. RSP makes it explicit rather than hidden. Availability mitigations — multi-signer Purpose Layer configurations requiring quorum approval for escalation — are architecturally feasible but introduce coordination latency. Competence mitigations — structured directive templates requiring the human to confirm specific diagnostic fields before issuing \texttt{RESOLVE} or \texttt{REJECT} — are within scope and reduce cognitive burden on the human anchor.

The harmful directive scenario warrants specific attention: an \texttt{ACK} that authorizes continued chain growth toward a harmful outcome based on incomplete diagnostic information or adversarial framing. A bounded \emph{dissent mechanism} — by which a node records a dissent annotation in the forensic ledger when it believes a directive is inconsistent with the chain's recorded intent, without the ability to override the directive — would strengthen accountability in Purpose Layer error or compromise scenarios. The dissent annotation creates a permanent, tamper-evident trail for post-hoc review. The mechanism must be strictly bounded: it must not create a pathway by which spawned agents refuse legitimate directives.

% ----------------------------------------------------------------------------%
\clearpage
\section{Conclusion}
\label{sec:rs-conclusion}

The Recursive Spawning Protocol makes a single architectural claim: accountability in multi-agent systems is achievable not through policy or convention, but through structure — specifically, through the immutable parent pointer, a cryptographically signed, Fact Layer-enforced reference from each child node to its parent, established at provisioning and unalterable by any actor after the fact.

This claim matters because the alternative — tracing agentic accountability through documentation, organizational policy, or operational convention — fails precisely under the conditions where accountability is most needed. When an agent spawns a child that spawns a child, when exception conditions propagate across multiple levels of delegation, when chain depth exceeds what any human can monitor in real time — under these conditions, convention-based accountability collapses. The delegation chain becomes an unanchored graph, and deeply spawned agents' actions become attributable to no one. RSP was designed for exactly these conditions.

The protocol's three correctness properties establish the operational meaning of the immutable parent pointer. \emph{Drift prevention} guarantees that every node in a spawn chain operates within a scope that is a descendant refinement of the root human's authorized intent — the chain cannot migrate purpose away from its human anchor through iterative scope expansion. \emph{Chain integrity} guarantees that chain topology and node configurations are tamper-evident after provisioning — no actor, including a compromised node, can alter parentage or configuration to obscure accountability. \emph{Termination} guarantees that the chain cannot grow indefinitely — the depth bound is enforced by the Fact Layer as a structural invariant, not applied as a policy convention.

These three properties are jointly enforced by the \bxthree{} Framework's three-layer architecture. The Purpose Layer defines spawn authorization policy and receives escalated exceptions, serving as the exclusive terminus of every accountability path. The Bounds Engine evaluates spawn necessity and enforces the depth constraint, serving as the constrained execution layer that cannot bypass the depth bound through operational urgency. The Fact Layer enforces the immutable parent pointer, maintains the append-only forensic ledger, and executes the pre-commit hooks that make all other constraints structurally enforceable. This mapping is not a design pattern — it is a structural necessity. Removing any layer causes the corresponding guarantee to collapse.

The immutable parent pointer is the architectural foundation of RSP precisely because it is the single structural element that cannot be circumvented, approximated, or weakened without breaking the chain of accountability entirely. Unlike depth limits, which can be bypassed through operational pressure or administrative override, the parent pointer is physically enforced by the Fact Layer at every spawn event. Unlike policy conventions, which depend on human compliance, the parent pointer is a property of the spawn event's structure — it exists in the forensic record whether or not any human monitors it. Unlike mutable delegation chains, which can be retroactively rewritten by any actor with sufficient access, the immutable parent pointer is append-only and hash-chained, making any tampering detectable at verification time.

This structural foundation produces an operational guarantee independent of chain depth, spawn frequency, or the operational urgency of the conditions that triggered spawning: every spawned agent is accountable to a named human, traceable through an immutable and verifiably complete parent pointer chain, and bounded by a Purpose Layer-defined depth that guarantees termination. The chain grows toward resolution capability. It does not grow away from accountability.

The limitations of the protocol — performance overhead from three-layer validation, scalability constraints from static depth bounds and the single-ledger architecture, and the compromised Purpose Layer problem — are the unavoidable cost of this guarantee. Each motivates a specific future work direction that preserves correctness guarantees while extending the protocol toward deployment scale and operational flexibility.

RSP does not resolve the problem of human accountability in AI systems generally. It resolves the structurally distinct problem of accountability in multi-agent spawn chains — the case in which agents spawn children and the delegation chain would otherwise collapse into an unanchored graph. For this problem, the immutable parent pointer is a definitive architectural answer, not a provisional mitigation.